Following \cite{lesnik2021synthetictopologyconstructivemetric}, with Andrej Bauer (and Claude), we explored an alternative notion of closed subsets which we will call \emph{weakly closed}.

\begin{definition}
  A proposition \(C\) is \notion{weakly closed}, if its negation is an open proposition.
\end{definition}

\begin{example}
  Any closed proposition is weakly closed. The proposition that an element of \(R\) is nilpotent is weakly closed.
\end{example}

The first point comes from the definitions in \cite{lesnik2021synthetictopologyconstructivemetric}[p. 47], where we choose ``T'' and ``\(\Sigma\)'' to be the type of propositions and compute that ``\(Z\)'' is the type of weakly closed propositions using \cite{lesnik2021synthetictopologyconstructivemetric}[Theorem 2.53].

\begin{remark}
  \begin{enumerate}[(a)]
  \item Let \(U\) be an open proposition, then \(\neg U\) is weakly closed.
  \item A proposition \(P\) is weakly closed if and only if \(P\to U\) is open for all open propositions \(U\).
  \item Weakly closed propositions are stable under finite disjunctions and conjuctions.
  \item A type \(X\) is compact, if and only if for all
    weakly closed \(C\subseteq X\), the proposition \((\exists x:X.x\in C)\) is weakly closed.
  \end{enumerate}
\end{remark}

\begin{proof}
  \begin{enumerate}[(a)]
  \item \(\neg\neg U=U\)
  \item Let \(U\) be open and \(P\) have the property that \(\neg P\) is open.
    Then we have:
    \begin{align*}
      (P\to U) = (P\to\neg\neg U) = (\neg(P\wedge \neg U)) = \neg\neg (\neg P\vee U) = (\neg P\vee U)
    \end{align*}
    \(\neg P \vee U\) is open, so \(P\to U\) is open.
  \item \(P\wedge Q\to U = P\to (Q\to U)\) is open for open \(U\) and weakly closed \(P,Q\).
    \(\neg(P\vee Q)= \neg P \wedge \neg Q\) is open since opens are closed under conjunction.
  \item \(\neg(\exists x:X.x\in C)=\forall x:X.\neg (x\in C)\).
  \end{enumerate}
\end{proof}
